% ЕРӨНХИЙ ДҮГНЭЛТ

\phantomsection
\addchaptertocentry{Ерөнхий дүгнэлт}
\label{Summary}

\section*{Ерөнхий дүгнэлт}

Энэхүү дипломын ажлын хүрээнд Монгол хэлний цахим орчин дахь хандлагыг
тодорхойлох, хортой агуулгыг бүтээлч шүүмжлэлээс ялган шүүх хиймэл оюунд
суурилсан NLP системийг боловсруулсан. Судалгааны гол үр дүнг дараах байдлаар
нэгтгэн дүгнэж болно.

\textbf{1. Хоёр шатлалт ангиллын архитектур.}
Уламжлалт гурван ангиллын (Positive / Neutral / Negative) сентимент шинжилгээг
цааш нь нарийвчлан, Negative агуулгыг Toxic болон Constructive гэсэн хоёр
дэд-ангилалд ялгах хоёр шатлалт (two-stage) архитектурыг зохион байгуулсан.
Энэхүү архитектур нь зөвхөн хортой агуулгыг устгах бус, бүтээлч шүүмжлэлийг
хамгаалах зорилгыг хангахад чиглэсэн.

\textbf{2. Өгөгдлийн бэлтгэл.}
news.mn болон Facebook платформуудаас Монгол кирилл бичгийн сэтгэгдлүүдийг
цуглуулж, кирилл шүүлтийн алхам, preprocessing pipeline, болон дөрвөн ангиллын
аннотацийн стратеги боловсруулсан. Кирилл бичгийн хязгаарлалт нь MN-BERT-ийн
токенчлолын тогтвортой байдлыг хангахад зайлшгүй шаардлагатай алхам болсон.

\textbf{3. MN-BERT загварын fine-tuning.}
Монгол хэлний урьдчилсан сургалттай MN-BERT загварыг Stage~1 (гурван ангилал)
болон Stage~2 (хоёр ангилал) тус тусдаа fine-tuning хийсэн. MN-BERT загвар нь
TF-IDF дээр суурилсан суурь загваруудтай (Logistic Regression, Naive Bayes,
SVM) харьцуулахад илүү сайн гүйцэтгэлтэй байсан.
% TODO: тодорхой тоон үр дүнг нэмнэ (жишээ: ``MN-BERT загвар Stage~1 дээр
% macro $F_1$=X.XX, Stage~2 дээр $F_1$=X.XX гүйцэтгэлтэй байсан'')

\textbf{4. Streamlit интерфейс.}
Хэрэглэгчийн интерфейстэй прототип системийг Streamlit framework ашиглан
хэрэгжүүлсэн бөгөөд хэрэглэгч текст оруулж ангиллын үр дүнг шууд харах,
batch шинжилгээ хийх боломжтой.

%-------------------------------------------------------------------------------
\section*{Цаашдын судалгааны чиглэл}

Дипломын ажлын хүрээнд хэрэгжүүлсэн системийг цаашид дараах чиглэлүүдээр
сайжруулах боломжтой:

\begin{enumerate}
  \item \textbf{Латин болон code-switched текстийн дэмжлэг:} Одоогийн систем
        зөвхөн кирилл бичгийн текстэд хамаарна. Цахим орчинд латин үсгийн
        бичвэр, монгол-англи холимог текст нэлээд түгээмэл тул эдгээрийг дэмжих
        нь хамрах хүрээг мэдэгдэхүйц өргөжүүлнэ.

  \item \textbf{Мультимодаль шинжилгээ:} Зураг, видео, emoji зэрэг текстэн бус
        элементүүдийг шинжилгээнд оруулах нь агуулгын илүү иж бүрэн ойлголтыг
        бий болгоно.

  \item \textbf{Илүү том dataset:} Цуглуулсан өгөгдлийн хэмжээг нэмэгдүүлж,
        олон платформ (X, TikTok, Instagram) хамруулах нь загварын ерөнхий
        чадварыг (generalization) сайжруулна.

  \item \textbf{Тайлбарлах чадвартай загвар (Explainable AI):} Загвар яагаад
        тухайн шийдвэр гаргаснаа тайлбарлах чадвар нэмэх нь хэрэглэгчийн
        итгэлцлийг нэмэгдүүлж, систем хариуцлагатай байх шаардлагыг хангана.

  \item \textbf{Бодит цагийн хяналт (real-time monitoring):} Системийг бодит
        цагийн агуулгын хяналтын платформ болгон хөгжүүлж, API хэлбэрээр бусад
        системтэй холбогдох боломжийг нээх.
\end{enumerate}
